% !TEX program = pdflatex
% !TEX options = -shell-escape -interaction=nonstopmode "%DOC%"
% Split wrapper: extracts the research.gov upload PDFs from main.pdf.
% Build main.pdf first, then either:
%   make split                                         (recommended)
%   pdflatex -shell-escape split.tex                   (writes both PDFs)
%   pdflatex -jobname=project-description split.tex   (pages 1-15)
%   pdflatex -jobname=references-cited split.tex      (page 16 to the end)
% Uses pdfTeX so the output is not rewritten by Preview/Quartz (rejected by research.gov).
% If the project description length changes, update \pdlast below.
\documentclass{article}
\usepackage{pdfpages}
\newcommand{\pdlast}{15}
\begin{document}
\ifnum\pdfstrcmp{\jobname}{project-description}=0
  \includepdf[pages=1-\pdlast]{main.pdf}
\else\ifnum\pdfstrcmp{\jobname}{references-cited}=0
  \includepdf[pages=\the\numexpr\pdlast+1\relax-]{main.pdf}
\else
  % Compiled directly as split.tex: generate both PDFs if shell-escape is on.
  \ifnum\pdfshellescape=1
    \immediate\write18{pdflatex -interaction=nonstopmode -jobname=project-description split.tex}
    \immediate\write18{pdflatex -interaction=nonstopmode -jobname=references-cited split.tex}
    \typeout{split.tex: wrote project-description.pdf and references-cited.pdf}
    \noindent Wrote \texttt{project-description.pdf} and \texttt{references-cited.pdf}.
  \else
    \typeout{split.tex: shell-escape is off; run `make split' instead}
    \noindent Shell-escape is off, so nothing was split.
    Run \texttt{make split} or \texttt{pdflatex -shell-escape split.tex}.
  \fi
\fi\fi
\end{document}
